% P033 only. Coordinates and incircle are calculated from the actual triangle.
% Arguments: A_x,A_y,C_x,scale,BC label,CA label,AB label,mode,caption.
% B=(0,0), C=(C_x,0). mode 0: circle; 1: areas; 2: congruence.
\providecommand{\PIncenterFigure}[9]{%
\begin{center}
\begin{tikzpicture}[scale=#4,font=\small,line join=round]
\coordinate (A) at (#1,#2); \coordinate (B) at (0,0); \coordinate (C) at (#3,0);
\pgfmathsetmacro{\pa}{#3}
\pgfmathsetmacro{\pb}{sqrt((#1-#3)^2+#2^2)}
\pgfmathsetmacro{\pc}{sqrt(#1^2+#2^2)}
\pgfmathsetmacro{\px}{(\pa*#1+\pc*#3)/(\pa+\pb+\pc)}
\pgfmathsetmacro{\pr}{\pa*#2/(\pa+\pb+\pc)}
\coordinate (I) at (\px,\pr);
\coordinate (D) at ($(B)!(I)!(C)$);
\coordinate (E) at ($(C)!(I)!(A)$);
\coordinate (F) at ($(A)!(I)!(B)$);
\ifnum#8=1
  \fill[blue!13] (I)--(B)--(C)--cycle;
  \fill[orange!20] (I)--(C)--(A)--cycle;
  \fill[green!17!white] (I)--(A)--(B)--cycle;
\fi
\ifnum#8=2
  \fill[orange!20] (A)--(E)--(I)--cycle;
  \fill[green!17!white] (A)--(F)--(I)--cycle;
  \draw[very thick,CProof] (A)--(I);
  \draw[CProof,dashed] (B)--(I);
  \pic[draw,CProof,angle radius=9mm] {angle=F--A--I};
  \pic[draw,CProof,angle radius=11mm] {angle=I--A--E};
\fi
\ifnum#8=0 \draw[CStatement,thick] (I) circle[radius=\pr];\fi
\draw[thick] (A)--node[pos=.30,left=5pt] {$#7$} (B)
 --node[midway,below=15pt] {$#5$} (C)
 --node[pos=.85,right=5pt] {$#6$} cycle;
\ifdim#1pt=0pt
  \pic[draw,angle radius=3mm] {right angle=C--B--A};
\fi
\ifdim#1pt=#3pt
  \pic[draw,angle radius=3mm] {right angle=A--C--B};
\fi
\ifnum#8=1
  \draw[gray,dashed] (I)--(A) (I)--(B) (I)--(C);
\fi
\draw[very thick,CStatement] (I)--node[midway,right=2pt,fill=white,inner sep=1pt] {$r$} (D);
\draw[very thick,CStatement] (I)--(E) (I)--(F);
\pic[draw,angle radius=2mm] {right angle=I--D--C};
\pic[draw,angle radius=2mm] {right angle=I--E--A};
\pic[draw,angle radius=2mm] {right angle=I--F--B};
\foreach \p/\where in {A/above,B/below left,C/below right,D/below,E/above right,F/above left}
 {\fill (\p) circle[radius=1.3pt]; \node[\where=3pt] at (\p) {$\p$};}
\fill[CStatement] (I) circle[radius=1.8pt];
\node[above=3pt,fill=white,inner sep=1pt] at (I) {$I$};
\node[anchor=north,align=center,text width=11cm] at (#3/2,-1.25/#4) {#9};
\end{tikzpicture}
\end{center}}

% A local angle isolates the two right triangles; no distant legend is needed.
\providecommand{\PAngleFigure}{%
\begin{center}
\begin{tikzpicture}[font=\small,line join=round]
\coordinate (A) at (0,0); \coordinate (E) at (4.5,2.598076);
\coordinate (F) at (4.5,-2.598076); \coordinate (I) at (6,0);
\fill[orange!18] (A)--(E)--(I)--cycle;
\fill[green!17!white] (A)--(F)--(I)--cycle;
\draw[thick] (A)--(5.5,3.175426) (A)--(5.5,-3.175426);
\draw[very thick,CProof] (A)--node[midway,above=3pt] {共同斜边 $AI$} (I);
\draw[very thick,CStatement] (E)--node[midway,right=3pt] {$IE$} (I)
 --node[midway,right=3pt] {$IF$} (F);
\pic[draw,angle radius=3mm] {right angle=A--E--I};
\pic[draw,angle radius=3mm] {right angle=I--F--A};
\pic[draw,angle radius=12mm] {angle=I--A--E};
\pic[draw,angle radius=14mm] {angle=F--A--I};
\foreach \p/\where in {A/left,E/above,F/below,I/right}
 {\fill (\p) circle[radius=1.5pt];\node[\where=4pt] at (\p) {$\p$};}
\node[align=center] at (3,-3.8) {两角相等、斜边共用 $\Longrightarrow$ 两直角三角形全等\\[3pt]
 对应的垂线段相等：$IE=IF$};
\end{tikzpicture}
\end{center}}

% P033: exact 13-14-15 triangle, A=(9,12), B=(0,0), C=(14,0).
% Optional values are labels only: x,y,z or the corresponding 7,8,6.
\providecommand{\PTangentFigure}[3]{%
\begin{center}
\begin{tikzpicture}[scale=.52,font=\small,line join=round]
\coordinate (A) at (9,12); \coordinate (B) at (0,0);
\coordinate (C) at (14,0); \coordinate (I) at (8,4);
\coordinate (D) at (8,0);
\coordinate (E) at ({152/13},{72/13});
\coordinate (F) at (4.8,6.4);
\draw[CStatement,thick,fill=blue!4] (I) circle[radius=4];
\draw[gray,dashed] (I)--(D) (I)--(E) (I)--(F);
\draw[CStatement,very thick] (A)--node[midway,left=5pt] {$#1$} (F)
 (A)--node[midway,right=5pt] {$#1$} (E);
\draw[CExample,very thick] (B)--node[midway,left=5pt] {$#2$} (F)
 (B)--node[midway,below=5pt] {$#2$} (D);
\draw[CProof,very thick] (C)--node[midway,right=5pt] {$#3$} (E)
 (C)--node[midway,below=5pt] {$#3$} (D);
\pic[draw,angle radius=2mm] {right angle=I--D--C};
\pic[draw,angle radius=2mm] {right angle=I--E--A};
\pic[draw,angle radius=2mm] {right angle=I--F--B};
\foreach \p/\where in {A/above,B/below left,C/below right,D/below,E/above right,F/above left}
 {\fill (\p) circle[radius=1.5pt];\node[\where=4pt] at (\p) {$\p$};}
\fill (I) circle[radius=1.6pt]; \node[above=4pt] at (I) {$I$};
\node[align=center,anchor=north,text width=11cm] at (7,-1.6)
 {同一顶点引出的两条切线段等长：每一对用同一种颜色标出。};
\end{tikzpicture}
\end{center}}

% P033: coordinates use capital X,Y to distinguish axes from tangent lengths x,y,z.
% 0 equilateral (side 6); 1 right (legs 3,4); 2 isosceles right (legs 4,4).
\providecommand{\PSpecialTriangle}[1]{%
\begin{center}
\begin{tikzpicture}[font=\small,line join=round,scale=1.05]
\ifnum#1=0
 \def\su{6}\def\sv{5.196152423}\def\sx{3}\def\si{3}\def\sr{1.732050808}
\else
 \def\sx{0}
 \ifnum#1=1
  \def\su{3}\def\sv{4}\def\si{1}\def\sr{1}
 \else
  \def\su{4}\def\sv{4}\def\si{1.171572875}\def\sr{1.171572875}
 \fi
\fi
\coordinate (B) at (0,0); \coordinate (C) at (\su,0);
\coordinate (A) at (\sx,\sv); \coordinate (I) at (\si,\sr);
\coordinate (D) at (\si,0); \coordinate (F) at (0,\sr);
\draw[gray,-{Stealth}] (-.7,0)--({\su+1},0) node[right] {$X$};
\draw[gray,-{Stealth}] (0,-.6)--(0,{\sv+1}) node[above] {$Y$};
\draw[CStatement,thick,fill=blue!5] (I) circle[radius=\sr];
\draw[thick] (A)--(B)--(C)--cycle;
\ifnum#1=0
 \draw[CProof,dashed] (A)--(D);
 \draw[CStatement,very thick] (I)--node[right=3pt] {$r$} (D);
 \node[below=6pt] at (D) {$\ell/2$};
 \node[below=6pt] at (C) {$\ell$};
 \node[left=5pt] at (0,\sv) {$h$};
 \draw[gray,dotted] (0,\sv)--(A);
 \node[left=5pt] at ($(A)!.48!(B)$) {$\ell$};
 \node[right=5pt] at ($(A)!.48!(C)$) {$\ell$};
 \node[above=5pt] at (A) {$A$};
 \node[above right=3pt] at (I) {$I$};
\else
 \draw[CStatement,very thick] (I)--(D) (I)--(F);
 \node[below=5pt] at (D) {$r$};
 \node[left=5pt] at (F) {$r$};
 \node[above left=4pt] at (I) {$I$};
 \pic[draw,angle radius=3mm] {right angle=C--B--A};
 \ifnum#1=1
  \node[below=6pt] at (C) {$u$};
  \node[left=5pt] at (A) {$v$};
  \node[right=5pt] at ($(A)!.32!(C)$) {$w$};
 \else
  \node[below=6pt] at (C) {$t$};
  \node[left=5pt] at (A) {$t$};
  \node[right=5pt] at ($(A)!.32!(C)$) {$\sqrt2t$};
  \draw[CProof,dashed] (B)--({\su/2},{\sv/2});
 \fi
 \node[above right=3pt] at (A) {$A$};
\fi
\node[below left=4pt] at (B) {$B=O$};
\node[above right=4pt] at (C) {$C$};
\foreach \p in {A,B,C,I} \fill (\p) circle[radius=1.5pt];
\end{tikzpicture}
\end{center}}

\begin{statementbox}
\textbf{先看图：内心到三边的垂线段，就是三个相等的半径。}
\PIncenterFigure{2}{4}{6}{1.05}{a}{b}{c}{0}{$ID=IE=IF=r$；距离沿垂线测量，不是连向顶点。}
\medskip
\textbf{\textcolor{CStatement}{条件和符号}}
\begin{description}[style=nextline,leftmargin=2.8em,labelsep=0.5em,itemsep=0.35em]
  \item[\textbf{条件1（三角形确实存在）：}] $\triangle ABC$ 的三个顶点不在同一直线上，
  $BC=a,\ CA=b,\ AB=c$，并且 $a,b,c>0$，任意两边之和大于第三边。
  \item[\textbf{条件2（内心和内切圆）：}] 三条\textbf{内角平分线}交于 $I$，
  过 $I$ 向三边 $BC,CA,AB$ 分别作垂线，垂足为 $D,E,F$。
  内切圆以 $I$ 为圆心，半径为 $r$。
  \item[\textbf{约定：}] 原三角形面积为 $S$，周长为 $L=a+b+c$，
  半周长为 $p=\dfrac{a+b+c}{2}$。
\end{description}

\medskip
\textbf{\textcolor{CStatement}{四个直接可用的结论}}
\begin{description}[style=nextline,leftmargin=2.8em,labelsep=0.5em,itemsep=0.45em]
  \item[\textbf{结论一（到三边距离相同）：}]
  \[
    ID=IE=IF=r.
  \]
  \textbf{注意：}这里的距离是\textbf{垂直距离}，不是 $IA,IB,IC$。
  \item[\textbf{结论二（半径、面积、周长互相换算）：}]
  \[
    \boxed{S=\frac12r(a+b+c)=rp=\frac12rL},\qquad
    \boxed{r=\frac{2S}{a+b+c}=\frac Sp}.
  \]
  \item[\textbf{结论三（三小块面积比等于对应底边比）：}]
  \[
    S_{\triangle IBC}:S_{\triangle ICA}:S_{\triangle IAB}=a:b:c.
  \]
\end{description}

\medskip
\par\noindent\begin{minipage}{\linewidth}
\textbf{\textcolor{CStatement}{结论四：已知三边，求各顶点到切点的距离}}

$D,E,F$ 也是内切圆与 $BC,CA,AB$ 的切点。沿三角形的边量取：
\[
 x=AE=AF,\qquad y=BF=BD,\qquad z=CD=CE.
\]
\PTangentFigure{x}{y}{z}
\end{minipage}\par
仍约定 $a=BC,b=CA,c=AB$，$p=(a+b+c)/2$，则
\[
\boxed{x=p-a=\frac{b+c-a}{2}},\qquad
\boxed{y=p-b=\frac{a+c-b}{2}},
\]
\[
\boxed{z=p-c=\frac{a+b-c}{2}}.
\]
\textbf{记法：}从哪个顶点出发，就用半周长减去该顶点的\textbf{对边长}。
三边满足三角形不等式，故 $x,y,z$ 均为正数。

\textbf{\textcolor{CStatement}{使用提示：}}
若只知道三边长度，可先求面积，再求半径；若已知半径和周长，可直接求面积；
若题目要求内心分出的三块面积比，甚至无需计算 $r$；
若只求切点分段，直接用半周长减对边，无需求面积或半径。
\end{statementbox}
